\documentclass[12pt,a4paper]{article}
\usepackage[UTF8,heading=true]{ctex}
\usepackage{geometry}
\usepackage{amsmath,amssymb}
\usepackage{enumitem,hyperref,fancyhdr,xcolor,setspace}
\usepackage{ragged2e}

\geometry{left=2.5cm,right=2.5cm,top=2.5cm,bottom=2.5cm}
\setlength{\headheight}{15pt}
\setstretch{1.5}
\setlength{\emergencystretch}{3em}
\hfuzz=20pt
\sloppy
\hypersetup{hidelinks}

\pagestyle{fancy}
\fancyhf{}
\fancyhead[C]{权利要求书}
\fancyfoot[C]{\thepage}
\renewcommand{\headrulewidth}{0pt}

\begin{document}

\begin{center}
    {\Huge\heiti 中国移动专利申请} \\[0.5em]
    {\Huge\heiti 权利要求书} \\[1em]
\end{center}

\noindent
发明名称：基于硅基光电集成的微型化QKD芯片封装与校准技术

\vspace{0.8em}

\begin{enumerate}[label=\arabic*.,leftmargin=2em,itemsep=0.6em]

\item
一种面向量子密钥分发QKD的硅基光子集成芯片封装结构，其特征在于，包括：承载基座、硅基光子QKD芯片、光纤或光波导耦合组件、电连接组件以及封装应力调控组件；其中，
所述硅基光子QKD芯片包括量子主光路与参考光路，所述参考光路与所述量子主光路在芯片版图上至少部分共模布设以共享温度场与封装应力场；
所述硅基光子QKD芯片还包括片上传感单元与片上执行单元，所述片上传感单元用于获得与温度漂移和/或应力漂移相关的传感信息，所述片上执行单元用于对所述量子主光路的相位、干涉平衡和/或耦合状态进行可控调节；
所述封装应力调控组件用于降低或重分配封装引入的残余应力并形成可预测的应力传递路径，以使所述片上传感单元与所述片上执行单元构成可闭环补偿的漂移通道。

\item
根据权利要求1所述的封装结构，其特征在于，所述参考光路包括参考波导、参考微环和/或参考马赫-曾德尔干涉器MZI；所述参考波导与所述量子主光路的关键干涉臂在几何长度、走线方向和邻近间距上满足预设共模约束，以提高对环境扰动的共模抑制能力。

\item
根据权利要求1所述的封装结构，其特征在于，所述片上传感单元包括热敏电阻、PN结温度传感器、微型热电堆、布拉格光栅传感器、压阻式应变计和/或基于环形谐振腔的应力敏感结构中的至少一种，用于分别输出温度观测量与应力观测量。

\item
根据权利要求1所述的封装结构，其特征在于，所述片上执行单元包括加热器阵列、电光相位调制器、可调衰减器和/或可调耦合器中的至少一种；所述加热器阵列沿所述量子主光路与所述参考光路的耦合漂移敏感区域分区布设，以支持多变量补偿。

\item
根据权利要求1所述的封装结构，其特征在于，所述电连接组件包括用于驱动所述片上执行单元的驱动电路和用于采集所述片上传感单元信号的采集电路，所述驱动电路和/或采集电路集成于本地控制芯片ASIC或控制微控制器MCU中。

\item
根据权利要求1所述的封装结构，其特征在于，所述封装应力调控组件包括对称固定点、低热膨胀系数基座材料、弹性粘接区和/或应力释放槽结构中的至少一种，以降低封装固化收缩与终端装配挤压导致的耦合失配与干涉偏置漂移。

\item
根据权利要求1所述的封装结构，其特征在于，所述参考光路在密钥分发时段内被降低功率、间歇注入或关闭，并在校准时段内被提升功率或切换到校准波长，以降低参考光对量子探测器的串扰污染风险。

\item
一种用于权利要求1至7任一项所述封装结构的片上自校准方法，其特征在于，包括：
获取片上传感单元输出的温度观测量与应力观测量，并获取参考光路输出的参考干涉响应或参考谐振响应；
基于所述温度观测量、所述应力观测量以及所述参考响应，估计量子主光路的漂移状态，并生成补偿控制量；
驱动片上执行单元执行补偿控制量，以维持量子主光路的相位设定、干涉可见度设定和/或耦合设定。

\item
根据权利要求8所述的自校准方法，其特征在于，所述估计漂移状态包括构建温度-应力联合漂移模型，并对所述漂移状态进行状态估计；所述状态估计采用递推估计、最小二乘估计、卡尔曼滤波和/或粒子滤波中的至少一种。

\item
根据权利要求8所述的自校准方法，其特征在于，所述自校准方法在密钥分发业务时序中插入静默时隙；在所述静默时隙内暂停或降载量子态发送并执行参考测量与补偿，以在不显著降低系统密钥吞吐的条件下实现在线校准。

\item
根据权利要求10所述的自校准方法，其特征在于，所述静默时隙的触发由以下至少一种条件决定：终端温度变化率超过阈值、终端姿态或加速度变化超过阈值、封装应力观测量突变超过阈值、量子误码率QBER或干涉可见度偏离阈值、参考谐振峰漂移超过阈值。

\item
根据权利要求8所述的自校准方法，其特征在于，所述补偿控制量同时包括针对温度漂移与应力漂移的联合补偿分量，并对不同漂移敏感区的加热器分区设置不同的权重或优先级，以降低热补偿引入的串扰与功耗。

\item
根据权利要求8所述的自校准方法，其特征在于，所述参考响应包括参考MZI的输出光强、参考微环的谐振波长或透射谱特征；所述方法进一步包括对参考光路的波长锁定，以使参考测量在预设的工作点附近进行。

\item
一种QKD模组或终端装置，其特征在于，包括：权利要求1至7任一项所述的封装结构、参考光源与量子光源、探测单元以及控制单元；
所述控制单元被配置为执行权利要求8至13任一项所述的自校准方法，并根据终端约束对校准频度、参考光功率和/或执行单元驱动策略进行自适应调整。

\item
根据权利要求14所述的QKD模组或终端装置，其特征在于，所述终端约束包括以下至少一种：电池电量或功耗预算、终端SoC发热状态、可用散热路径、体积或厚度限制、业务时延限制、量子信道占空比限制；所述控制单元在满足所述终端约束的前提下维持预设的密钥生成率下限或QBER上限。

\item
根据权利要求14所述的QKD模组或终端装置，其特征在于，所述控制单元还被配置为与对端节点协同，在密钥分发协议的帧结构内对静默时隙进行对齐、指示或确认，以避免校准过程对协议同步与安全参数评估产生不利影响。

\item
一种电子设备，其特征在于，包括处理器、存储器以及权利要求14至16任一项所述的QKD模组或终端装置，所述存储器中存储有可由所述处理器执行的程序，所述程序在执行时用于实现权利要求8至13任一项所述的自校准方法。

\item
一种计算机可读存储介质，其上存储有计算机程序，所述计算机程序在被处理器执行时用于实现权利要求8至13任一项所述的自校准方法。

\end{enumerate}

\end{document}

